\section{Diskussion}
\textcolor{red}{• von der Statistik zur Fragestellung
• Interpretation der statistischen Ergebnisse im Kontext der ursprünglichen Fragestellung
sind die Ergebnisse plausibel/überraschend?
was bedeuten die Ergebnisse
• statistische Diskussion der Simulationsergebnisse}

\subsection{Modellwahl}

Es wurden vier Kandidatenmodelle auf Basis biologischer Überlegungen erstellt und das beste Modell mittels BIC
ermittelt. Das Kriterium wägt zwischen Modelllgüte und Komplexität ab und identifiziert so das am besten
passende Modell. Anschließend wurde das beste Modell validiert, indem die Residuen, der LOOCV, k-fache Kreuzvalidierung
und der Shapiro Test, Breusch-Pagan Test, sowie die Cooks Distanz betrachtet wurden.
\\
Das finale Modell \texttt{lm1} enthält nur drei Variablen mit dem Intercept, kann damit aber die Daten
gut schätzen. Es treten Über- und Unterschätzungen nur im geringen Maß auf siehe Abbildung \ref{fig:Residuen}.
Das zeigt die gewählten Prädikatoren erfassen noch nicht alle wichtigen Zusammenhänge. Mögliche Ursachen könnten fehlende
Einflussgrößen oder Interaktionen sein. Jedoch können auf Grundlage der 38 Datenpunkte keine weiteren
Zusammenhänge zuverlässig ermittelt werden, wodurch das sparsamste Modell am statistisch stabilsten ist.
\\
Die Modelle sind von den Validierungsergebnissen nah beieinander, jedoch zeigt das finale Modell den
niedrigsten BIC und vergleichbare Validierungswerte. Dies zeigt das die Temperatur und die Dicke der Streuschicht wichtige Faktoren sind, die die Daten beschreiben.
Die Bodentemperatur gibt die kinetische Energie für alle biochemischen Reaktionen in den Zellen der Bodenorganismen.
Dadurch ist sie ein wichtiger zeitlicher Treiber. Die Temperatur in 15 cm tiefe hat eine hohe Korrelationen (0,43),
dies ist biologisch sinnvoll, da die Temperaturen in tieferen Schichten stabiler sind und damit besser die
Wurzel und Mikrobenpopulation repräsentieren. Durch den Quadratischen Term wird abgebildet, dass sich die
biologische Aktivität bei steigenden Temperaturen überproportional beschleunigt und ein Otpimum besitzt.
\\
Die Streuschichtdicke repräsentiert die Quantität des frischen organischen
Materials, das den Mikroorganismen als primäre Energiequelle für die heterotrophe Respiration dient.
Während die Temperatur primär zeitliche Schwankungen erklärt, ist die Dicke der Streu in Wäldern ein
entscheidender Faktor für die räumlichen Muster der Bodenatmung. Der quadratische Term deutet darauf hin,
dass die Menge des Substrats die Atmung nicht einfach nur linear erhöht, sondern dass z. B. eine sehr
dicke Streuschicht die mikrobielle Aktivität durch bessere Feuchtigkeitskonservierung oder ein massives
Nahrungsüberangebot überproportional fördert. Die linearen Terme wurden weggelassen, das verletzt zwar
das übliche statistische Hierarchie-Prinzip, jedoch zwang die geringe Stichprobengröße zu dieser radikalen
Reduktion, um die Modellstabilität (BIC) zu maximieren.
\\
Interessant ist das \texttt{lmoi} nicht in dem Modell vorkommt, obwohl es die stärkste Korrelation zur Zielgröße
aufweist und auch bei \cite{diss} im Modell genutzt wird. Das könnte daran liegen das die Feuchte der
Streuschicht mit der Dicke der Streuschicht mit 0,37 korriliert. Weitere Gründe könnten sein, dass das
Signal bei unserem kleinen Stichprobenumfang nicht stark genug war, um sich vom Rauschen abzuheben oder
die Schätzung mit dem BIC zu instabil ist.
\\
Das finale Modell berücksichtigt nur physikalisch-chemische Faktoren, die biotischen Faktoren wie die
Feinwurzelbiomasse fehlen. Es wird also nur die heterotrophe und nicht die autotrophe Atmung abgebildet,
oderbwohl die Wurzelatmung bis zu 90\% der Bodenatmung ausmachen kann \cite{diss}.
Dies könnte daran liegen, dass \texttt{temp.15} mit der Trockenmasse der Wurzeln in 0-5 cm tiefe mit 0,37
korriliert. Und das zu wenige Daten vorliegen um weitere biologische Effekte zu identifizieren. Zusätzlich könnte
es auch sein, dass die Temperatur auch die Aktivität der Wurzeln abbildet.
\\
Vorallem bei wenigen Daten ist die Gefahr des Overfittings sehr groß, wobei streng das
Parsimonie-Prinzip angewendet werden muss. Insgesamt zeigen die Ergebnisse das ein Modell mit wenigen
Parametern schon gut Daten beschreiben kann.

\subsection{Validierung}

\subsection{Simulation}
Das Ziel der Simulationsstudie ist die falsche Anwendung des t-Tests bei Vorwärtsselektion von
Prädiktoren. Die Ergebnisse der Simulationsstudie in \ref{ErgebnisSim} zeigt das die Verteilungen
der T-max Werte rechtsschief sind. Zu erwarten bei einem T-Test ist eine Normalverteilung also
keinerlei Schiefe. Die nicht Normalität in den Ergebnissen ist zu erklären dadurch das es sich
hier nicht um die Verteilung eines einzelnen T-Test handelt sondern um eine Verteilung der
Maxima vieler T-Tests. Das führt auch zu erheblich höheren Werten zusätzlich zu der
Rechtsschiefe. Da wir weiterhin den klassischen Wert $\alpha = 0.05$ als kritischen Wert
verwenden, werden viel mehr Prädiktoren aufgenommen als im vorher festgelegten besten Modell
enthalten sind. In diesem Fall wurden sieben Prädiktoren aufgenommen obwohl das Modell was
verwendet wurde um die Pseudowerte dieser Simulation zu erzeugen nur zwei Prädiktoren beinhaltet. 

Diese erhöte Aufnahme führt zu einer einer generell extrem hohen Ablehnung der Nullhypothese und
damit erheblich erhöhten Fehlerwahrscheinlichkeit erster Art. Das könnte man beheben in dem man
den kritischen Wert anpasst zum Beispiel durch eine Bonferroni Korrektur oder einfach eine
Simulation des max T. 

Selektions-Bias ist ein weiterer großer Schwachpunkt dieser Simulation. In der ersten Runde
zeigt die Aufnahme eines Prädiktors zwangsweise die beste Reduktion der Residuen im Vergleich
zu den anderen Prädiktoren. In diesem Fall $temp.15^2$. Da der T-wert dieses Prädiktors am
höchsten ist und höher als der kritische Wert wird dieser aufgenommen aber es kann noch andere
Prädiktoren geben die über dem kritischen Wert liegen und vielleicht mit Interaktionen mit
anderen Prädiktoren eine besserere Erklärungskraft haben. Diese werden hier aber nicht
aufgenommen und potenziell in kommenden Selektionsrunden auch nicht.

Zudem kommt noch das Problem von Overfitting, es ist durchaus möglich das die Simulation das
best möglichste Modell aufstellt für die Erklärung der Daten auf denen es trainiert ist
allerdings kann es sein das es nicht besonders gute Vorhersagen gegeben anderer Daten treffen
wird. Ein großer Faktor hierfür ist das immer die Variable aufgenommen wird die die Residuen
am stärksten reduziert selbst wenn diese Reduktion sehr minimal ist und potenziel in anderen
Bedingungen große Auswirkungen hat.